\section{Relatividad general y cosmología homogénea}
\label{sec:rg}

Esta sección fija el lenguaje y las ecuaciones de la RG que se usan en todo el informe. Es material estándar \etq{E}; se incluye porque la gravedad unimodular se define \emph{por contraste} con la RG, y para entenderla hace falta tener claro qué es exactamente lo que cambia.

\subsection{La acción de Einstein--Hilbert y sus ecuaciones}

La RG describe la gravedad mediante una métrica $g_{ab}$ sobre el espacio-tiempo, cuya dinámica sale de extremizar la acción de Einstein--Hilbert más la acción de la materia,
\begin{equation}
S=\frac{1}{2\kap}\int \dd^4x\,\sqrt{-g}\,R+S_m[g_{ab},\Psi],
\label{eq:EH}
\end{equation}
donde $R$ es el escalar de Ricci, $g=\det g_{ab}$ y $\Psi$ representa los campos de materia. Se define el tensor energía-impulso de la materia como la respuesta de $S_m$ a una variación de la métrica,
\begin{equation}
T_{ab}\equiv-\frac{2}{\sqrt{-g}}\,\frac{\delta S_m}{\delta g^{ab}}.
\label{eq:Tab}
\end{equation}
Con la variación estándar $\delta(\sqrt{-g}R)=\sqrt{-g}\,G_{ab}\,\delta g^{ab}$ (salvo un término de borde), donde $G_{ab}\equiv R_{ab}-\tfrac12 R\,g_{ab}$ es el tensor de Einstein, anular la variación respecto de $g^{ab}$ da las ecuaciones de Einstein,
\begin{equation}
G_{ab}=\kap\,T_{ab}.
\label{eq:Einstein}
\end{equation}
Son diez ecuaciones (el tensor es simétrico) para las diez componentes de la métrica. El lado izquierdo describe la geometría y el derecho la materia.

Hay un hecho que será central más adelante. El tensor de Einstein satisface la identidad geométrica de Bianchi, $\nabla^a G_{ab}=0$, que vale para \emph{cualquier} métrica. Aplicada a \eqref{eq:Einstein} implica
\begin{equation}
\nabla^a T_{ab}=0,
\label{eq:conservacion}
\end{equation}
es decir, la energía-impulso de la materia se conserva. En RG esto no es una hipótesis adicional: es una consecuencia de las ecuaciones de campo (y, de otro modo, de la invarianza de $S_m$ frente a difeomorfismos cuando la materia satisface sus ecuaciones de movimiento). Retener esta observación es la clave para entender por qué la gravedad unimodular es distinta.

\subsection{La constante cosmológica}

Se puede agregar a la acción un término $-2\Lambda$ dentro del integrando de \eqref{eq:EH}, de modo que $S=\frac{1}{2\kap}\int\sqrt{-g}\,(R-2\Lambda)+S_m$. Las ecuaciones de campo pasan a ser
\begin{equation}
G_{ab}+\Lambda\,g_{ab}=\kap\,T_{ab}.
\label{eq:EinsteinLambda}
\end{equation}
Como $\Lambda$ es una constante y $\nabla_a g^{ab}=0$, la conservación \eqref{eq:conservacion} sigue valiendo. El término $\Lambda g_{ab}$ equivale a un fluido de densidad $\rho_\Lambda=\Lambda/\kap$ y presión $p_\Lambda=-\rho_\Lambda$, que es la ecuación de estado del vacío. En RG, $\Lambda$ es un \emph{parámetro de la teoría}: un número que hay que medir. Y aquí aparece la dificultad que motiva a la gravedad unimodular: cualquier contribución de energía del vacío de los campos de materia tiene el mismo efecto que $\Lambda$, de modo que el valor efectivo que gravita es la suma de una constante desnuda y de esas contribuciones, y el valor observado es inmensamente menor que las estimaciones naturales (véase la introducción de \cite{leon22} para la discusión de ese desacuerdo, y \cite{weinberg89} como referencia clásica, que no se leyó para este trabajo).

\subsection{El universo homogéneo e isótropo}

A escalas cosmológicas el universo es, con muy buena aproximación, homogéneo e isótropo. Con secciones espaciales planas, su métrica es la de FLRW,
\begin{equation}
\dd s^2=-\dd t^2+a(t)^2\,\delta_{ij}\,\dd x^i\dd x^j,
\label{eq:FLRW}
\end{equation}
donde $a(t)$ es el factor de escala y $t$ el tiempo cósmico (el de un observador comóvil). El parámetro de Hubble es $H\equiv\dot a/a$. Los símbolos de Christoffel no nulos son $\Gamma^0_{ij}=a\dot a\,\delta_{ij}$ y $\Gamma^i_{0j}=H\delta^i_j$; las componentes del tensor de Ricci son $R_{00}=-3\ddot a/a$ y $R_{ij}=(a\ddot a+2\dot a^2)\delta_{ij}=a^2(\dot H+3H^2)\delta_{ij}$; el escalar de Ricci es $R=6(\dot H+2H^2)$, y el tensor de Einstein en forma mixta resulta
\begin{equation}
G^0{}_0=-3H^2,\qquad G^i{}_j=-(2\dot H+3H^2)\,\delta^i_j.
\label{eq:Gfondo}
\end{equation}
La materia homogénea e isótropa se describe como un fluido perfecto, con $T^0{}_0=-\rho$ y $T^i{}_j=p\,\delta^i_j$, siendo $\rho(t)$ la densidad de energía y $p(t)$ la presión. Las componentes $00$ e $ij$ de \eqref{eq:Einstein} son las ecuaciones de Friedmann,
\begin{equation}
3H^2=\kap\rho,\qquad 2\dot H+3H^2=-\kap p.
\label{eq:Friedmann}
\end{equation}
Restándolas se obtiene una relación que resultará útil por su forma,
\begin{equation}
\dot H=-\frac{\kap}{2}\,(\rho+p),
\label{eq:Hdot}
\end{equation}
que dice que $H$ solo decrece si $\rho+p>0$, y que $\dot H$ \emph{solo ve la combinación $\rho+p$}. Derivando la primera ecuación de \eqref{eq:Friedmann} y usando \eqref{eq:Hdot} se obtiene la ecuación de continuidad, que expresa la conservación \eqref{eq:conservacion} en este contexto,
\begin{equation}
\dot\rho+3H(\rho+p)=0.
\label{eq:continuidad}
\end{equation}
Si además se suma un término $\Lambda$, las ecuaciones cambian a $3H^2=\kap\rho+\Lambda$ y $2\dot H+3H^2=-\kap p+\Lambda$; \eqref{eq:Hdot} y \eqref{eq:continuidad} no cambian.

De \eqref{eq:Friedmann} sale también la ecuación de aceleración,
\begin{equation}
\frac{\ddot a}{a}=\dot H+H^2=-\frac{\kap}{6}\,(\rho+3p),
\label{eq:aceleracion}
\end{equation}
de modo que la expansión se acelera ($\ddot a>0$) si y solo si $\rho+3p<0$. Con una ecuación de estado $p=w\rho$, la continuidad da $\rho\propto a^{-3(1+w)}$, y la aceleración exige $w<-1/3$. La materia ordinaria ($w=0$) y la radiación ($w=1/3$) desaceleran; solo algo parecido a un fluido de energía de vacío ($w\approx-1$) acelera. Esa es la observación de la que parte la inflación.

Conviene registrar dos variables auxiliares que se usarán de manera recurrente. El tiempo conforme $\eta$ se define por $\dd\eta=\dd t/a$, y con él la métrica \eqref{eq:FLRW} queda $\dd s^2=a^2(-\dd\eta^2+\dd\mathbf x^2)$; las cantidades derivadas respecto de $\eta$ se anotan con prima, y $\mathcal H=a'/a=aH$. Y el número de e-folds $N=\ln(a/a_i)$ mide el tiempo en unidades de la expansión: $\dd N=H\,\dd t$. En estas variables, $a''/a=\mathcal H'+\mathcal H^2=a^2(2H^2+\dot H)$, cuyo uso se verá en las secciones~\ref{sec:tensores} y~\ref{sec:numerico}.
